\documentclass[11pt, a4paper]{report}

% ========== Common Packages ==========
\usepackage{fontspec}
\usepackage{kotex}
\usepackage{amsmath, amssymb, amsthm}
\usepackage{mathtools}
\usepackage{bm}
\usepackage{graphicx}
\usepackage{booktabs}
\usepackage{array}
\usepackage{tabularx}
\usepackage{longtable}
\usepackage{hyperref}
\usepackage{cleveref}
\usepackage{geometry}
\usepackage{fancyhdr}
\usepackage{titlesec}
\usepackage{enumitem}
\usepackage{xcolor}
\usepackage{tcolorbox}
\usepackage{algorithm}
\usepackage{algorithmic}
\usepackage{tikz}
\usetikzlibrary{arrows.meta, positioning, calc, decorations.pathreplacing, shapes.geometric, fit, patterns, angles, quotes, trees}
\usepackage{pgfplots}
\pgfplotsset{compat=1.18}
\usepackage{caption}
\usepackage{subcaption}
\usepackage{float}

\geometry{margin=2.5cm, top=3cm, bottom=3cm}

\usepackage{multirow}
\usepackage{siunitx}
% ========== Colors ==========
% Common
\definecolor{defblue}{RGB}{30, 80, 150}
\definecolor{thmgreen}{RGB}{20, 110, 60}
\definecolor{noteorange}{RGB}{200, 100, 0}
\definecolor{lightblue}{RGB}{230, 240, 255}
\definecolor{lightgreen}{RGB}{230, 255, 235}
\definecolor{lightorange}{RGB}{255, 245, 230}
\definecolor{lightgray}{RGB}{245, 245, 245}
% Extended
\definecolor{lightred}{RGB}{255, 235, 235}
\definecolor{darkred}{RGB}{180, 30, 30}
\definecolor{biogreen}{RGB}{10, 130, 80}
\definecolor{cvpurple}{RGB}{100, 40, 150}
\definecolor{injuryred}{RGB}{200, 40, 40}
\definecolor{codegray}{RGB}{240, 240, 245}
\definecolor{pipeblue}{RGB}{25, 100, 180}
\definecolor{constraintred}{RGB}{200, 50, 50}
\definecolor{termblack}{RGB}{30, 30, 30}
\definecolor{goldcolor}{RGB}{180, 140, 20}

% ========== Theorem-like environments ==========
\theoremstyle{definition}
\newtheorem{definition}{Definition}[chapter]
\newtheorem{example}{Example}[chapter]

\theoremstyle{plain}
\newtheorem{theorem}{Theorem}[chapter]
\newtheorem{proposition}{Proposition}[chapter]

\theoremstyle{remark}
\newtheorem{remark}{Remark}[chapter]

% ========== tcolorbox environments ==========
% --- Common (all parts) ---
\newtcolorbox{keyidea}[1][]{
  colback=lightblue,
  colframe=defblue,
  fonttitle=\bfseries,
  title={Key Idea},
  #1
}

\newtcolorbox{importantnote}[1][]{
  colback=lightorange,
  colframe=noteorange,
  fonttitle=\bfseries,
  title={Important},
  #1
}

\newtcolorbox{summary}[1][]{
  colback=lightgreen,
  colframe=thmgreen,
  fonttitle=\bfseries,
  title={Summary},
  #1
}

% --- Warning/Error ---
\newtcolorbox{warningbox}[1][]{
  colback=lightred,
  colframe=darkred,
  fonttitle=\bfseries,
  title={Warning},
  #1
}

% --- Domain: Biomechanics & Clinical ---
\newtcolorbox{clinicalnote}[1][]{
  colback=lightred,
  colframe=darkred,
  fonttitle=\bfseries,
  title={Clinical Note},
  #1
}

\newtcolorbox{biomechbox}[1][]{
  colback=lightgreen,
  colframe=biogreen,
  fonttitle=\bfseries,
  title={Biomechanics},
  #1
}

\newtcolorbox{injurybox}[1][]{
  colback={injuryred!8},
  colframe=injuryred,
  fonttitle=\bfseries,
  title={Injury Risk},
  #1
}

% --- Domain: Computer Vision ---
\newtcolorbox{cvbox}[1][]{
  colback={cvpurple!5},
  colframe=cvpurple,
  fonttitle=\bfseries,
  title={Computer Vision},
  #1
}

% --- Pipeline & Setup ---
\newtcolorbox{setupbox}[1][]{
  colback=lightgray,
  colframe=gray!60,
  fonttitle=\bfseries,
  title={Setup},
  #1
}

\newtcolorbox{pipelinebox}[1][]{
  colback=pipeblue!5,
  colframe=pipeblue,
  fonttitle=\bfseries,
  title={Pipeline Step},
  #1
}

\newtcolorbox{codebox}[1][]{
  colback=codegray,
  colframe=gray!70,
  fonttitle=\bfseries\ttfamily,
  title={Code},
  #1
}

\newtcolorbox{termbox}[1][]{
  colback=termblack!5,
  colframe=termblack!60,
  fonttitle=\bfseries\ttfamily,
  title={Terminal},
  #1
}

% --- Research ---
\newtcolorbox{hypothesis}[1][]{
  colback=goldcolor!8,
  colframe=goldcolor,
  fonttitle=\bfseries,
  title={Hypothesis},
  #1
}

\newtcolorbox{resultbox}[1][]{
  colback=thmgreen!8,
  colframe=thmgreen,
  fonttitle=\bfseries,
  title={Expected Result},
  #1
}

% --- Constraints & Solutions ---
\newtcolorbox{constraintbox}[1][]{
  colback=constraintred!8,
  colframe=constraintred,
  fonttitle=\bfseries,
  title={Constraint},
  #1
}

\newtcolorbox{solutionbox}[1][]{
  colback=thmgreen!8,
  colframe=thmgreen,
  fonttitle=\bfseries,
  title={Solution},
  #1
}

\newtcolorbox{databox}[1][]{
  colback=pipeblue!5,
  colframe=pipeblue,
  fonttitle=\bfseries,
  title={Data Spec},
  #1
}

% --- Progress/Status ---
\newtcolorbox{successbox}[1][]{
  colback=lightgreen,
  colframe=thmgreen,
  fonttitle=\bfseries,
  title={Success},
  #1
}

\newtcolorbox{nextbox}[1][]{
  colback=lightorange,
  colframe=noteorange,
  fonttitle=\bfseries,
  title={Next Step},
  #1
}

% --- Fitting guide ---
\newtcolorbox{failbox}[1][]{
  colback=lightred,
  colframe=darkred,
  fonttitle=\bfseries,
  title={Failure Pattern},
  #1
}

\newtcolorbox{fixbox}[1][]{
  colback=lightgreen,
  colframe=thmgreen,
  fonttitle=\bfseries,
  title={Fix},
  #1
}

\newtcolorbox{lessonbox}[1][]{
  colback=lightorange,
  colframe=noteorange,
  fonttitle=\bfseries,
  title={Lesson Learned},
  #1
}

% --- Specification ---
\newtcolorbox{specbox}[1][]{
  colback=cvpurple!5,
  colframe=cvpurple,
  fonttitle=\bfseries,
  title={Specification},
  #1
}

\newtcolorbox{checklistbox}[1][]{
  colback=lightgreen,
  colframe=biogreen,
  fonttitle=\bfseries,
  title={Checklist},
  #1
}

% ========== Custom math commands ==========
% Vectors (lowercase)
\newcommand{\vx}{\mathbf{x}}
\newcommand{\vd}{\mathbf{d}}
\newcommand{\vo}{\mathbf{o}}
\newcommand{\vc}{\mathbf{c}}
\newcommand{\vr}{\mathbf{r}}
\newcommand{\vp}{\mathbf{p}}
\newcommand{\vv}{\mathbf{v}}
\newcommand{\vn}{\mathbf{n}}
\newcommand{\vu}{\mathbf{u}}
\newcommand{\vt}{\mathbf{t}}
\newcommand{\ve}{\mathbf{e}}
\newcommand{\vq}{\mathbf{q}}
% Vectors (uppercase)
\newcommand{\vF}{\mathbf{F}}
\newcommand{\vM}{\mathbf{M}}
\newcommand{\va}{\mathbf{a}}
\newcommand{\vg}{\mathbf{g}}
\newcommand{\vX}{\mathbf{X}}
% Bold Greek
\newcommand{\vmu}{\bm{\mu}}
\newcommand{\vbeta}{\bm{\beta}}
\newcommand{\vtheta}{\bm{\theta}}
\newcommand{\vpsi}{\bm{\psi}}
% Matrices
\newcommand{\mSigma}{\bm{\Sigma}}
\newcommand{\mR}{\mathbf{R}}
\newcommand{\mS}{\mathbf{S}}
\newcommand{\mW}{\mathbf{W}}
\newcommand{\mJ}{\mathbf{J}}
\newcommand{\mG}{\mathbf{G}}
\newcommand{\mT}{\mathbf{T}}
\newcommand{\mI}{\mathbf{I}}
\newcommand{\mB}{\mathbf{B}}
\newcommand{\mK}{\mathbf{K}}
\newcommand{\mP}{\mathbf{P}}
\newcommand{\mF}{\mathbf{F}}
\newcommand{\mE}{\mathbf{E}}
\newcommand{\mA}{\mathbf{A}}
\newcommand{\mH}{\mathbf{H}}
% Sets and groups
\newcommand{\RR}{\mathbb{R}}
\newcommand{\SO}{\mathrm{SO}}
% Units
\newcommand{\degps}{\,\text{deg/s}}
\newcommand{\Nm}{\ensuremath{\,\text{N}\cdot\text{m}}}
\newcommand{\BW}{\,\text{BW}}

% ========== Listings ==========
\usepackage{listings}

\lstset{
  basicstyle=\ttfamily\small,
  keywordstyle=\color{defblue}\bfseries,
  commentstyle=\color{thmgreen}\itshape,
  stringstyle=\color{noteorange},
  backgroundcolor=\color{codegray},
  frame=single,
  rulecolor=\color{gray!40},
  breaklines=true,
  showstringspaces=false,
  numbers=left,
  numberstyle=\tiny\color{gray},
  tabsize=4,
  captionpos=b
}

% ========== Header/Footer ==========
\pagestyle{fancy}
\fancyhf{}
\fancyhead[L]{\leftmark}
\fancyhead[R]{\thepage}
\renewcommand{\headrulewidth}{0.4pt}

% ========== Title formatting ==========
\titleformat{\chapter}[display]
  {\normalfont\huge\bfseries\color{defblue}}
  {\chaptertitlename\ \thechapter}{20pt}{\Huge}

\titleformat{\section}
  {\normalfont\Large\bfseries\color{defblue!80!black}}
  {\thesection}{1em}{}

\titleformat{\subsection}
  {\normalfont\large\bfseries\color{defblue!60!black}}
  {\thesubsection}{1em}{}


\begin{document}

% ==================== TITLE PAGE ====================
\begin{titlepage}
\centering
\vspace*{1.5cm}
{\Large\color{gray} Part 4}
\vspace{0.5cm}

{\Huge\bfseries\color{defblue} 다시점 촬영 데이터\\[0.3cm] 분석 파이프라인}

\vspace{1cm}

{\LARGE\color{defblue!70!black} 데이터 수집부터 3DGS 인체 재구성,\\생체역학 지표 산출까지}

\vspace{2cm}

\begin{tikzpicture}[
  scale=0.7,
  block/.style={rectangle, rounded corners, draw=defblue, fill=lightblue,
    minimum width=2.8cm, minimum height=0.9cm, font=\small\bfseries, text=defblue!80!black},
  arrow/.style={-Stealth, thick, defblue!70}
]
% Pipeline blocks
\node[block] (data) at (0,0) {Data Collection};
\node[block] (calib) at (4,0) {Calibration};
\node[block] (pose2d) at (8,0) {2D Pose};
\node[block] (smpl) at (0,-2) {SMPL-X Fit};
\node[block] (gart) at (4,-2) {GART / 3DGS};
\node[block] (bio) at (8,-2) {Biomechanics};
\node[block] (valid) at (4,-4) {Validation};

\draw[arrow] (data) -- (calib);
\draw[arrow] (calib) -- (pose2d);
\draw[arrow] (pose2d) -- (smpl);
\draw[arrow] (smpl) -- (gart);
\draw[arrow] (gart) -- (bio);
\draw[arrow] (bio) -- (valid);
\draw[arrow] (smpl) -- (valid);
\end{tikzpicture}

\vspace{2cm}
{\large 컴퓨터 공학 $\times$ 스포츠의학}
\vspace{0.5cm}

{\large 2026}

\vspace{\fill}
{\small\color{gray} 본 교재는 Part 1 (NeRF$\rightarrow$GART), Part 2 (3DGS+동작분석), Part 3 (투구 실험설계)의 후속편으로,\\
촬영된 데이터를 실제로 처리$\cdot$분석하는 전체 파이프라인을 단계별로 다룹니다.}
\end{titlepage}

\tableofcontents
\newpage

% ==================== PREFACE ====================
\chapter*{서문}
\addcontentsline{toc}{chapter}{서문}

Part 3까지는 ``무엇을 촬영할 것인가''와 ``어떤 장비를 어떻게 배치할 것인가''에 집중했다.
실험 설계가 끝나면 자연스럽게 떠오르는 질문은 이것이다:
\textbf{``촬영이 끝난 뒤, 이 수백 GB의 영상과 모션캡처 파일을 어떻게 처리하여 의미 있는 생체역학 지표로 바꿀 것인가?''}

Part 4는 이 질문에 대한 답이다. 7대의 카메라에서 나온 동기화 영상과 Vicon 마커 데이터를 입력으로 받아,
\begin{enumerate}[nosep]
  \item 카메라 보정과 좌표계 정합,
  \item 2D/3D 자세 추정,
  \item SMPL-X 파라메트릭 인체 모델 피팅,
  \item 3D Gaussian Splatting(GART) 기반 포토리얼리스틱 인체 재구성,
  \item 관절 각도$\cdot$토크$\cdot$힘 등 생체역학 지표 산출,
  \item Vicon gold-standard와의 정량적 검증
\end{enumerate}
까지 이르는 \textbf{end-to-end 데이터 분석 파이프라인}을 단계별로 설명한다.

각 챕터는 이론적 배경을 먼저 제시한 뒤, 실제로 사용하는 명령어와 코드 스니펫을 함께 보여 주어
독자가 직접 재현할 수 있도록 구성하였다. 특히 고속 투구 동작이라는 난이도 높은 시나리오를 기준으로 설명하되,
일반적인 보행 분석이나 다른 스포츠 동작에도 동일한 파이프라인을 적용할 수 있음을 강조한다.

\begin{importantnote}
본 교재의 코드 예시는 Python 3.10+, PyTorch 2.x, CUDA 11.8+ 환경을 기준으로 한다.
운영체제는 Ubuntu 22.04를 권장하며, macOS에서도 대부분의 전처리 단계를 수행할 수 있다.
\end{importantnote}

\newpage

% ==================== CHAPTER 1 ====================
\chapter{데이터 수집과 전처리}

\section{다시점 영상 데이터 구조}

촬영이 완료되면 가장 먼저 해야 할 일은 데이터를 체계적으로 정리하는 것이다.
일관된 디렉토리 구조는 이후 모든 자동화 스크립트의 전제 조건이 된다.

\begin{pipelinebox}[title={권장 디렉토리 레이아웃}]
\begin{lstlisting}[language=bash, numbers=none]
project/
  subject_01/
    set_001/
      device1_synced.MOV    # Camera 1
      device2_synced.MOV    # Camera 2
      ...
      device7_synced.MOV    # Camera 7
      _synced_meta.json     # Sync metadata
      sync_offsets.json     # Per-device time offsets
    set_002/
      ...
  calibration/
    intrinsics/             # Per-camera intrinsic params
    extrinsics/             # Multi-camera extrinsic params
    charuco_images/         # Calibration board images
\end{lstlisting}
\end{pipelinebox}

\begin{keyidea}
파일명에 \texttt{device\{N\}\_synced} 규칙을 사용하면 카메라 인덱스를 파싱하기 쉽고,
\texttt{\_synced} 접미사로 동기화 완료 여부를 즉시 확인할 수 있다.
\end{keyidea}

\subsection{네이밍 컨벤션}

다시점 촬영 시스템에서는 수십$\sim$수백 개의 파일이 생성되므로,
명확한 네이밍 규칙이 필수적이다.

\begin{itemize}[nosep]
  \item \textbf{피험자 ID}: \texttt{subject\_XX} (2자리 정수)
  \item \textbf{세트 번호}: \texttt{set\_XXX} (3자리 정수, 시행 순서)
  \item \textbf{카메라 ID}: \texttt{deviceN} ($N = 1, \ldots, 7$)
  \item \textbf{동기화 상태}: \texttt{\_synced} (동기화 완료) / \texttt{\_raw} (원본)
\end{itemize}

\subsection{메타데이터 파일}

\texttt{\_synced\_meta.json}에는 각 카메라 영상의 원본 경로, 프레임 레이트, 해상도,
녹화 시작$\cdot$종료 타임스탬프가 기록된다. \texttt{sync\_offsets.json}에는 각 디바이스 간의
시간 오프셋(밀리초 단위)이 저장된다.

\begin{codebox}[title={sync\_offsets.json 예시}]
\begin{lstlisting}[language=Python, numbers=none]
{
  "device1": 0.0,      # Reference camera
  "device2": -3.2,     # 3.2ms earlier
  "device3": 1.5,      # 1.5ms later
  ...
}
\end{lstlisting}
\end{codebox}

\section{프레임 동기화 검증}

다시점 촬영에서 가장 중요한 전처리 단계는 \textbf{시간 동기화(temporal synchronization)}이다.
카메라 간 프레임이 정확히 같은 시점을 가리키지 않으면, 삼각측량 결과에 체계적 오차가 발생한다.

\subsection{동기화 방법}

세 가지 대표적 동기화 방법을 비교한다.

\begin{table}[H]
\centering
\caption{다시점 동기화 방법 비교}
\begin{tabularx}{\textwidth}{lXcc}
\toprule
\textbf{방법} & \textbf{원리} & \textbf{정밀도} & \textbf{난이도} \\
\midrule
하드웨어 젠록 & 공통 클럭 신호로 셔터 동기 & $<0.1$\,ms & 높음 \\
LED 플래시 & 촬영 시작 시 LED 점멸, 프레임에서 검출 & $\sim$1 프레임 & 낮음 \\
오디오 싱크 & 클랩보드/비프음의 파형 상관 분석 & $\sim$1\,ms & 중간 \\
\bottomrule
\end{tabularx}
\end{table}

\subsection{소프트웨어 기반 동기화 검증}

하드웨어 동기화를 사용하더라도 사후 검증이 필수적이다.
대표적 방법은 빠르게 움직이는 객체(예: 투구 중 공)의 위치를 각 카메라에서 추적한 뒤,
에피폴라 제약(epipolar constraint)을 만족하는지 확인하는 것이다.

\begin{definition}[에피폴라 제약]
두 카메라 $i, j$의 기본 행렬(fundamental matrix)을 $\mathbf{F}_{ij}$라 하면,
대응점 $\vx_i, \vx_j$에 대해
\begin{equation}
  \vx_j^\top \mathbf{F}_{ij}\, \vx_i = 0
\end{equation}
이 성립해야 한다. 동기화 오차가 존재하면 이 값이 0에서 벗어난다.
\end{definition}

\begin{remark}
240\,Hz 촬영에서 1 프레임의 동기화 오차는 약 4.17\,ms에 해당한다.
투구 중 손 끝의 속도가 $\sim$30\,m/s라면, 1 프레임 오차는 약 12.5\,cm의 위치 오차를 야기한다.
따라서 고속 동작에서는 서브프레임 정밀도의 동기화가 필수적이다.
\end{remark}

\section{영상 전처리}

\subsection{디베이어링(Debayering)}

RAW 영상을 촬영한 경우, Bayer 패턴에서 풀 컬러 이미지로의 변환(디베이어링)이 필요하다.
일반적인 소비자 카메라(MOV/MP4)의 경우 카메라 내부에서 이미 처리되어 이 단계를 건너뛸 수 있다.

\subsection{렌즈 왜곡 보정}

광각 렌즈를 사용하면 방사형(radial) 및 접선형(tangential) 왜곡이 발생한다.
OpenCV의 왜곡 모델은 다음과 같다:

\begin{equation}
\begin{aligned}
  x_{\text{dist}} &= x(1 + k_1 r^2 + k_2 r^4 + k_3 r^6) + 2p_1 xy + p_2(r^2 + 2x^2) \\
  y_{\text{dist}} &= y(1 + k_1 r^2 + k_2 r^4 + k_3 r^6) + p_1(r^2 + 2y^2) + 2p_2 xy
\end{aligned}
\end{equation}
여기서 $r^2 = x^2 + y^2$, $(k_1, k_2, k_3)$는 방사형 왜곡 계수,
$(p_1, p_2)$는 접선형 왜곡 계수이다.

\begin{codebox}[title={OpenCV 왜곡 보정}]
\begin{lstlisting}[language=Python]
import cv2
import numpy as np

# Load calibration
K = np.load("intrinsics/cam1_K.npy")     # 3x3
dist = np.load("intrinsics/cam1_dist.npy")  # [k1,k2,p1,p2,k3]

# Undistort
map1, map2 = cv2.initUndistortRectifyMap(
    K, dist, None, K, (W, H), cv2.CV_32FC1
)
undistorted = cv2.remap(frame, map1, map2, cv2.INTER_LINEAR)
\end{lstlisting}
\end{codebox}

\subsection{색 보정과 화이트밸런스}

다시점 시스템에서 모든 카메라의 색 응답이 동일해야 뷰 간 일관된 학습이 가능하다.
촬영 전에 ColorChecker 차트를 촬영하여 카메라별 색 변환 행렬 $\mathbf{C}_i \in \RR^{3 \times 3}$을 추정한다:
\begin{equation}
  \mathbf{I}_{\text{corrected}} = \mathbf{C}_i \, \mathbf{I}_{\text{raw}}
\end{equation}

\section{Vicon 마커 데이터 전처리}

\subsection{갭 필링(Gap Filling)}

마커가 일시적으로 가려지면 궤적에 빈 구간(gap)이 생�다. Vicon Nexus에서는 다음 방법을 제공한다:

\begin{itemize}[nosep]
  \item \textbf{Spline Fill}: 3차 스플라인 보간 --- 짧은 갭($<$10 프레임)에 적합
  \item \textbf{Pattern Fill}: 인접 마커의 움직임 패턴을 참조 --- 긴 갭에 유용
  \item \textbf{Rigid Body Fill}: 강체 가정 하에 다른 마커로부터 복원
\end{itemize}

\subsection{라벨링}

자동 라벨링 후 반드시 수동 검증이 필요하다. 특히 고속 동작에서는 마커 교차(swap)가 빈번하게 발생하므로,
스켈레톤 모델을 오버레이하여 프레임별로 확인한다.

\subsection{필터링}

마커 궤적에 포함된 고주파 노이즈를 제거하기 위해 저역통과 필터를 적용한다.
투구 동작에서의 일반적 차단 주파수는 12--25\,Hz이다 (잔차 분석 기반 선택).

\begin{biomechbox}
\textbf{잔차 분석(Residual Analysis)} \cite{winter2009}: 다양한 차단 주파수 $f_c$에 대해
필터링된 신호와 원본의 잔차 RMS를 계산한다. 잔차가 직선적으로 증가하기 시작하는 지점이
최적 $f_c$이다. 이는 신호와 노이즈의 파워 스펙트럼이 분리되는 주파수에 해당한다.
\end{biomechbox}


% ==================== CHAPTER 2 ====================
\chapter{카메라 보정과 COLMAP}

\section{개별 카메라 내부 보정}

카메라 내부 파라미터(intrinsic parameters)는 3D 점을 2D 이미지 좌표로 매핑하는
투영 모델의 핵심이다.

\begin{definition}[카메라 내부 행렬]
핀홀 카메라의 내부 행렬은 다음과 같다:
\begin{equation}
  \mK = \begin{pmatrix} f_x & 0 & c_x \\ 0 & f_y & c_y \\ 0 & 0 & 1 \end{pmatrix}
\end{equation}
여기서 $(f_x, f_y)$는 초점 거리(픽셀 단위), $(c_x, c_y)$는 주점(principal point)이다.
\end{definition}

\subsection{체커보드 보정}

Zhang의 방법 \cite{zhang2000}을 사용하여 체커보드 패턴으로 내부 파라미터와 왜곡 계수를 동시에 추정한다.

\begin{codebox}[title={체커보드 보정 코드}]
\begin{lstlisting}[language=Python]
import cv2
import glob

BOARD = (9, 6)  # inner corners
objp = np.zeros((BOARD[0]*BOARD[1], 3), np.float32)
objp[:, :2] = np.mgrid[0:BOARD[0], 0:BOARD[1]].T.reshape(-1, 2)
objp *= 30.0  # 30mm square size

obj_pts, img_pts = [], []
for path in sorted(glob.glob("calib_images/cam1_*.png")):
    img = cv2.imread(path, cv2.IMREAD_GRAYSCALE)
    ret, corners = cv2.findChessboardCorners(img, BOARD)
    if ret:
        corners = cv2.cornerSubPix(img, corners, (11,11), (-1,-1),
            (cv2.TERM_CRITERIA_EPS+cv2.TERM_CRITERIA_COUNT, 30, 0.001))
        obj_pts.append(objp)
        img_pts.append(corners)

ret, K, dist, rvecs, tvecs = cv2.calibrateCamera(
    obj_pts, img_pts, img.shape[::-1], None, None
)
print(f"RMS reprojection error: {ret:.4f} px")
\end{lstlisting}
\end{codebox}

\begin{importantnote}
재투영 오차(RMS)가 0.5 픽셀을 초과하면 보정 이미지를 재검토해야 한다.
체커보드가 다양한 각도와 위치에서 촬영되었는지, 이미지 전체 영역을 커버하는지 확인한다.
\end{importantnote}

\subsection{ChArUco 보드}

ChArUco 보드는 체커보드와 ArUco 마커를 결합한 것으로, 부분적으로 가려져도 코너를 검출할 수 있어
다시점 환경에서 특히 유용하다.

\section{COLMAP SfM/MVS 파이프라인}

COLMAP은 Structure from Motion(SfM)과 Multi-View Stereo(MVS)를 수행하는 표준 도구로,
3D Gaussian Splatting의 초기 점군 생성에 핵심적으로 사용된다.

\begin{figure}[H]
\centering
\begin{tikzpicture}[
  node distance=1.2cm,
  block/.style={rectangle, rounded corners, draw=defblue, fill=lightblue,
    minimum width=3.5cm, minimum height=0.8cm, font=\small, text=defblue!80!black, align=center},
  arrow/.style={-Stealth, thick, defblue!60}
]
\node[block] (feat) {Feature Extraction\\(SIFT / SuperPoint)};
\node[block, below=of feat] (match) {Feature Matching\\(Exhaustive / Vocab Tree)};
\node[block, below=of match] (sfm) {Incremental SfM\\(Bundle Adjustment)};
\node[block, below=of sfm] (dense) {Dense Reconstruction\\(PatchMatch MVS)};
\node[block, below=of dense] (fuse) {Point Cloud Fusion\\(Depth Map Merging)};

\draw[arrow] (feat) -- (match);
\draw[arrow] (match) -- (sfm);
\draw[arrow] (sfm) -- (dense);
\draw[arrow] (dense) -- (fuse);

% Outputs
\node[right=2cm of sfm, font=\small\color{gray}] (out1) {cameras.bin, images.bin, points3D.bin};
\node[right=2cm of fuse, font=\small\color{gray}] (out2) {fused.ply};

\draw[-Stealth, gray, dashed] (sfm.east) -- (out1.west);
\draw[-Stealth, gray, dashed] (fuse.east) -- (out2.west);
\end{tikzpicture}
\caption{COLMAP 파이프라인 흐름도}
\end{figure}

\begin{codebox}[title={COLMAP CLI 사용법}]
\begin{lstlisting}[language=bash]
# 1. Feature extraction
colmap feature_extractor \
    --database_path db.db \
    --image_path images/ \
    --ImageReader.camera_model OPENCV \
    --ImageReader.single_camera_per_folder 1

# 2. Feature matching
colmap exhaustive_matcher --database_path db.db

# 3. Sparse reconstruction (SfM)
mkdir sparse
colmap mapper \
    --database_path db.db \
    --image_path images/ \
    --output_path sparse/

# 4. Dense reconstruction (optional for 3DGS)
colmap image_undistorter \
    --image_path images/ \
    --input_path sparse/0 \
    --output_path dense/
\end{lstlisting}
\end{codebox}

\subsection{Bundle Adjustment}

Bundle Adjustment(BA)는 모든 카메라 자세와 3D 점 위치를 동시에 최적화하는 비선형 최소제곱 문제이다:

\begin{equation}
  \min_{\{\mR_i, \vt_i\}, \{\vp_j\}} \sum_{i,j} \rho\!\left(\left\| \pi(\mR_i, \vt_i, \mK_i, \vp_j) - \vx_{ij} \right\|^2\right)
\end{equation}
여기서 $\pi$는 투영 함수, $\rho$는 로버스트 손실 함수(Cauchy/Huber), $\vx_{ij}$는 카메라 $i$에서의 점 $j$의 관측이다.

\section{다시점 외부 보정}

\subsection{동시 보정 vs 순차 보정}

\begin{table}[H]
\centering
\caption{외부 보정 전략 비교}
\begin{tabularx}{\textwidth}{lXX}
\toprule
\textbf{전략} & \textbf{장점} & \textbf{단점} \\
\midrule
동시(Simultaneous) & 모든 카메라 쌍의 관계를 한 번에 최적화, 전역 일관성 보장 & 대규모 시스템에서 계산 비용 높음 \\
순차(Sequential) & 카메라를 하나씩 추가하며 점진적 보정, 구현 용이 & 오차 누적(drift) 가능 \\
\bottomrule
\end{tabularx}
\end{table}

\subsection{알려진 내부 파라미터 활용}

사전에 보정된 내부 파라미터를 COLMAP에 고정(fix)하면 외부 보정의 정확도가 향상된다:

\begin{codebox}[title={내부 파라미터 고정}]
\begin{lstlisting}[language=bash]
colmap mapper \
    --Mapper.ba_refine_focal_length 0 \
    --Mapper.ba_refine_principal_point 0 \
    --Mapper.ba_refine_extra_params 0
\end{lstlisting}
\end{codebox}

\section{Vicon--RGB 좌표계 정합}

Vicon 시스템과 카메라 시스템은 서로 다른 좌표계를 사용하므로, 두 좌표계 간의 강체 변환(rigid transformation)을 추정해야 한다.

\begin{definition}[Procrustes 정합]
$N$개의 대응점 $\{\va_i\} \subset \RR^3$ (Vicon)과 $\{\mathbf{b}_i\} \subset \RR^3$ (카메라)이 주어졌을 때,
최적의 회전 $\mR^*$, 이동 $\vt^*$, 스케일 $s^*$를 구한다:
\begin{equation}
  \mR^*, \vt^*, s^* = \arg\min_{\mR, \vt, s} \sum_{i=1}^{N} \left\| \mathbf{b}_i - (s\mR\va_i + \vt) \right\|^2
\end{equation}
\end{definition}

실용적으로, Vicon 마커 중 최소 4개를 RGB 이미지에서도 식별 가능한 위치에 부착하여 공통 랜드마크로 활용한다.
반사 마커는 적외선에서만 밝게 보이므로, 별도의 컬러 마커나 체커보드 코너를 공통점으로 사용하는 것이 일반적이다.

\section{보정 품질 평가}

\subsection{재투영 오차}

카메라 보정의 가장 기본적인 품질 지표는 재투영 오차이다:
\begin{equation}
  e_{\text{reproj}} = \frac{1}{NM} \sum_{i=1}^{N}\sum_{j=1}^{M} \left\| \pi(\mR_i, \vt_i, \mK_i, \vp_j) - \vx_{ij} \right\|_2
\end{equation}

\begin{table}[H]
\centering
\caption{재투영 오차 기준}
\begin{tabular}{lc}
\toprule
\textbf{수준} & \textbf{재투영 오차 (px)} \\
\midrule
우수 & $< 0.3$ \\
양호 & $0.3 \sim 0.5$ \\
주의 필요 & $0.5 \sim 1.0$ \\
재보정 필요 & $> 1.0$ \\
\bottomrule
\end{tabular}
\end{table}

\subsection{삼각측량 정밀도}

두 카메라 이상에서 관측된 점의 3D 위치를 삼각측량으로 복원한 뒤,
복원 정밀도를 평가한다. 특히 기저선(baseline) 대비 깊이 비율이 중요하다:

\begin{equation}
  \sigma_z \propto \frac{z^2}{f \cdot B}
\end{equation}
여기서 $z$는 깊이, $f$는 초점 거리, $B$는 기저선 길이이다.
기저선이 짧으면 깊이 추정의 불확실성이 급격히 증가한다.


% ==================== CHAPTER 3 ====================
\chapter{2D 자세 추정과 다시점 삼각측량}

\section{프레임별 2D 자세 추정}

각 카메라 영상의 모든 프레임에서 2D 관절 위치를 추정한다.
현재 최고 성능의 모델은 ViTPose \cite{xu2022vitpose}와 RTMPose \cite{jiang2023rtmpose}이다.

\begin{cvbox}
\textbf{ViTPose}는 Vision Transformer 기반 자세 추정 모델로,
사전학습된 MAE(Masked Autoencoder) 백본 위에 간단한 디코더를 올린 구조이다.
COCO 벤치마크에서 AP 75.8 (ViTPose-H)을 달성하며,
큰 모델 크기 대비 높은 정확도를 제공한다.

\textbf{RTMPose}는 실시간 추론에 특화된 경량 모델로,
SimCC(Simplified Coordinate Classification) 기반 좌표 예측을 사용한다.
240fps 영상의 실시간 처리가 가능하다.
\end{cvbox}

\begin{codebox}[title={MMPose를 이용한 배치 추론}]
\begin{lstlisting}[language=Python]
from mmpose.apis import MMPoseInferencer

inferencer = MMPoseInferencer(
    pose2d='td-hm_ViTPose-huge_8xb64-210e_coco-256x192'
)

# Process all frames from one camera
results = inferencer(
    'frames/cam1/',
    show=False,
    out_dir='pose2d/cam1/',
    return_datasamples=True
)
\end{lstlisting}
\end{codebox}

\section{고속 영상의 특수 문제}

\subsection{모션 블러}

240\,Hz 이상에서도 투구 릴리스 순간의 손$\cdot$공은 심각한 모션 블러를 보일 수 있다.
셔터 속도를 $1/1000$\,s 이하로 설정하여 블러를 최소화해야 한다.

\begin{importantnote}
셔터 속도를 높이면 조리개를 더 열거나 ISO를 높여야 하므로
피사계 심도(DoF)가 얕아지거나 노이즈가 증가한다.
240\,Hz + $1/1000$\,s 조합에서는 충분한 조명(>2000 lux)이 필수적이다.
\end{importantnote}

\subsection{반복 질감(Repetitive Texture)}

실내 체육관의 균일한 바닥이나 단색 벽면은 특징점 매칭을 어렵게 만든다.
해결 방법:
\begin{itemize}[nosep]
  \item 바닥에 비반복적 패턴(예: 랜덤 도트) 부착
  \item SuperPoint/SuperGlue 같은 학습 기반 특징점 사용
  \item 매칭 시 에피폴라 제약을 적극 활용하여 잘못된 매칭 제거
\end{itemize}

\section{다시점 삼각측량}

$C$개의 카메라에서 관측된 2D 관절 위치 $\{\vx_c\}_{c=1}^{C}$로부터
3D 위치 $\vp \in \RR^3$를 복원한다.

\subsection{DLT(Direct Linear Transform)}

각 카메라의 투영 행렬을 $\mP_c = \mK_c[\mR_c | \vt_c] \in \RR^{3 \times 4}$라 하면,
관측 $\vx_c = (u_c, v_c)$에 대해 다음 선형 방정식을 세울 수 있다:

\begin{equation}
\begin{pmatrix}
u_c \, \vp_c^{3\top} - \vp_c^{1\top} \\
v_c \, \vp_c^{3\top} - \vp_c^{2\top}
\end{pmatrix}
\begin{pmatrix} X \\ Y \\ Z \\ 1 \end{pmatrix} = \mathbf{0}
\end{equation}
여기서 $\vp_c^{k\top}$는 $\mP_c$의 $k$번째 행이다.
모든 카메라의 방정식을 쌓아 $\mathbf{A}\vp_h = \mathbf{0}$ 형태로 만들고,
SVD로 풀면 최소제곱 해를 얻는다.

\begin{codebox}[title={DLT 삼각측량 구현}]
\begin{lstlisting}[language=Python]
def triangulate_dlt(proj_matrices, points_2d):
    """
    proj_matrices: list of (3,4) projection matrices
    points_2d: list of (2,) 2D observations
    returns: (3,) 3D point
    """
    n = len(proj_matrices)
    A = np.zeros((2*n, 4))
    for i, (P, x) in enumerate(zip(proj_matrices, points_2d)):
        A[2*i]   = x[0]*P[2] - P[0]
        A[2*i+1] = x[1]*P[2] - P[1]
    _, _, Vt = np.linalg.svd(A)
    X = Vt[-1]
    return X[:3] / X[3]
\end{lstlisting}
\end{codebox}

\subsection{RANSAC 기반 이상치 제거}

일부 카메라에서 2D 검출이 실패하거나 부정확할 수 있다.
RANSAC을 적용하여 이상치 뷰를 제거한다:

\begin{algorithm}[H]
\caption{RANSAC 삼각측량}
\begin{algorithmic}[1]
\REQUIRE 카메라 투영 행렬 $\{\mP_c\}$, 2D 관측 $\{\vx_c\}$, 임계값 $\tau$
\FOR{$k = 1$ to $K_{\max}$}
  \STATE 2개의 카메라를 랜덤 선택
  \STATE DLT로 3D 점 $\vp$ 추정
  \STATE 모든 카메라에 재투영: $\hat{\vx}_c = \pi(\mP_c, \vp)$
  \STATE 인라이어 집합: $\mathcal{I} = \{c : \|\hat{\vx}_c - \vx_c\| < \tau\}$
  \IF{$|\mathcal{I}|$ > 최대 인라이어 수}
    \STATE 최적 해 갱신
  \ENDIF
\ENDFOR
\STATE 최적 인라이어 집합으로 DLT 재수행
\RETURN $\vp^*$
\end{algorithmic}
\end{algorithm}

\section{시계열 필터링}

삼각측량으로 얻은 3D 관절 궤적에는 검출 노이즈로 인한 고주파 변동이 포함된다.
이를 제거하기 위해 시계열 필터를 적용한다.

\subsection{Butterworth 저역통과 필터}

$n$차 Butterworth 필터의 주파수 응답은:
\begin{equation}
  |H(f)|^2 = \frac{1}{1 + (f / f_c)^{2n}}
\end{equation}
여기서 $f_c$는 차단 주파수, $n$은 필터 차수이다.
위상 지연을 제거하기 위해 \textbf{양방향 필터링(zero-phase filtering)}을 적용한다.

\begin{codebox}[title={양방향 Butterworth 필터}]
\begin{lstlisting}[language=Python]
from scipy.signal import butter, filtfilt

def butterworth_filter(data, cutoff, fs, order=4):
    nyq = 0.5 * fs
    b, a = butter(order, cutoff / nyq, btype='low')
    return filtfilt(b, a, data, axis=0)

# 240 Hz data, 15 Hz cutoff
joints_filtered = butterworth_filter(joints_3d, 15, 240)
\end{lstlisting}
\end{codebox}

\subsection{Savitzky--Golay 필터}

다항식 피팅 기반의 Savitzky--Golay 필터는 신호의 피크를 잘 보존하는 장점이 있다:

\begin{equation}
  \hat{y}_i = \sum_{m=-M}^{M} c_m \, y_{i+m}
\end{equation}
여기서 $c_m$은 $2M+1$ 크기 윈도우 내에서 $p$차 다항식 피팅으로 결정되는 계수이다.

\begin{keyidea}
Butterworth는 주파수 영역에서, Savitzky--Golay는 시간 영역에서 동작한다.
투구처럼 충격적(impulsive) 신호가 포함된 경우 Savitzky--Golay가 피크 보존에 유리하고,
주기적 보행 동작에서는 Butterworth가 더 깔끔한 결과를 준다.
\end{keyidea}


% ==================== CHAPTER 4 ====================
\chapter{다시점 SMPL-X 피팅}

\section{EasyMocap 파이프라인 상세}

EasyMocap \cite{easymocap}은 다시점 영상에서 SMPL/SMPL-X 파라미터를 추정하는 오픈소스 프레임워크이다.
입력으로 카메라 보정 파라미터와 2D 자세 검출 결과를 받아, 최적화 기반으로 인체 모델을 피팅한다.

\begin{pipelinebox}[title={EasyMocap 처리 흐름}]
\begin{enumerate}[nosep]
  \item 카메라 파라미터 로드 (내부 + 외부)
  \item 프레임별 2D 관절 검출 결과 로드
  \item 다시점 삼각측량으로 초기 3D 관절 추정
  \item SMPL-X 모델 초기화 (평균 체형 + 중립 자세)
  \item 단계적 최적화: 전역 자세 $\rightarrow$ 체형 $\rightarrow$ 관절 자세 $\rightarrow$ 손/얼굴
  \item 결과 저장: 프레임별 SMPL-X 파라미터 $(\vbeta, \vtheta, \vt)$
\end{enumerate}
\end{pipelinebox}

\section{다시점 목적 함수와 최적화}

SMPL-X 피팅의 핵심은 다시점 2D 관측과 3D 모델을 일치시키는 목적 함수의 설계이다.

\begin{equation}
\mathcal{L}_{\text{total}} = \lambda_{\text{2D}} \mathcal{L}_{\text{2D}} + \lambda_{\text{3D}} \mathcal{L}_{\text{3D}} + \lambda_{\beta} \mathcal{L}_{\beta} + \lambda_{\theta} \mathcal{L}_{\theta} + \lambda_{\text{smooth}} \mathcal{L}_{\text{smooth}}
\end{equation}

\subsection{2D 재투영 손실}

$C$개 카메라의 2D 관측과 SMPL-X 관절의 투영을 비교한다:
\begin{equation}
  \mathcal{L}_{\text{2D}} = \sum_{c=1}^{C} \sum_{j=1}^{J} w_{cj} \left\| \pi_c\!\big(J_j(\vbeta, \vtheta)\big) - \hat{\vx}_{cj} \right\|^2
\end{equation}
여기서 $J_j(\vbeta, \vtheta)$는 SMPL-X의 $j$번째 관절의 3D 위치,
$\hat{\vx}_{cj}$는 카메라 $c$에서의 2D 검출, $w_{cj}$는 검출 신뢰도이다.

\subsection{3D 관절 손실}

삼각측량으로 얻은 3D 관절과의 차이를 최소화한다:
\begin{equation}
  \mathcal{L}_{\text{3D}} = \sum_{j=1}^{J} \left\| J_j(\vbeta, \vtheta) - \hat{\vp}_j \right\|^2
\end{equation}

\subsection{사전 확률 정규화}

체형과 자세의 과도한 변형을 방지한다:
\begin{equation}
  \mathcal{L}_{\beta} = \|\vbeta\|^2, \qquad
  \mathcal{L}_{\theta} = \|\vtheta\|^2
\end{equation}

\begin{importantnote}
$\mathcal{L}_{\theta}$에 단순 $L_2$ 정규화를 사용하면 모든 관절이 중립 자세(T-pose)로
당겨진다. 실제로는 VPoser \cite{pavlakos2019} 같은 학습된 자세 사전(prior)을 사용하여
해부학적으로 타당한 자세 공간 내에서 최적화하는 것이 바람직하다.
\end{importantnote}

\section{시간적 일관성 정규화}

프레임별 독립 피팅은 시간적 떨림(jitter)을 유발한다.
인접 프레임 간의 파라미터 변화를 제한하는 정규화를 추가한다:

\begin{equation}
  \mathcal{L}_{\text{smooth}} = \sum_{t=2}^{T} \left( \|\vtheta_t - \vtheta_{t-1}\|^2 + \gamma \|\vtheta_t - 2\vtheta_{t-1} + \vtheta_{t-2}\|^2 \right)
\end{equation}

첫째 항은 속도 정규화(1차 유한차분), 둘째 항은 가속도 정규화(2차 유한차분)이다.
$\gamma$는 일반적으로 0.5--2.0 사이로 설정한다.

\begin{keyidea}
가속도 정규화는 등속 운동은 허용하면서 급격한 방향 전환만 억제한다.
투구의 가속 구간(arm cocking)에서는 $\gamma$를 낮추고,
팔로-스루에서는 높이는 적응형 전략이 효과적이다.
\end{keyidea}

\section{피팅 결과 시각화 및 디버깅}

\subsection{체크리스트}

피팅 결과를 검증하기 위한 체계적 체크리스트:

\begin{enumerate}[nosep]
  \item \textbf{실루엣 오버레이}: 각 카메라 뷰에서 SMPL-X 메시를 원본 이미지에 오버레이
  \item \textbf{관절 재투영}: 3D 관절을 2D로 투영하여 검출 결과와 비교
  \item \textbf{체형 일관성}: 모든 프레임에서 $\vbeta$가 일정한지 확인 (같은 피험자이므로)
  \item \textbf{관절 범위}: 해부학적으로 불가능한 관절 각도가 없는지 확인
  \item \textbf{침투 검사}: 자기 교차(self-intersection)가 없는지 확인
\end{enumerate}

\begin{codebox}[title={피팅 결과 시각화}]
\begin{lstlisting}[language=Python]
import trimesh
import pyrender

# Load SMPL-X result
body_model = smplx.create(model_path, model_type='smplx')
output = body_model(betas=betas, body_pose=pose,
                    global_orient=orient, transl=transl)
vertices = output.vertices.detach().numpy()[0]
faces = body_model.faces

# Render overlay
mesh = trimesh.Trimesh(vertices, faces)
scene = pyrender.Scene()
scene.add(pyrender.Mesh.from_trimesh(mesh))
# ... add camera, render, blend with original image
\end{lstlisting}
\end{codebox}

\subsection{흔한 피팅 실패와 해결}

\begin{table}[H]
\centering
\caption{SMPL-X 피팅 실패 유형과 해결 방법}
\begin{tabularx}{\textwidth}{lXX}
\toprule
\textbf{증상} & \textbf{원인} & \textbf{해결} \\
\midrule
메시가 뒤집어짐 & 전역 회전 초기화 실패 & 초기 자세를 수동으로 설정 \\
팔다리 관통 & 자기 교차 정규화 부재 & $\mathcal{L}_{\text{penetration}}$ 추가 \\
떨림(jitter) & 시간 정규화 가중치 부족 & $\lambda_{\text{smooth}}$ 증가 \\
체형 왜곡 & 뷰 수 부족 & 카메라 추가 또는 $\lambda_\beta$ 증가 \\
\bottomrule
\end{tabularx}
\end{table}


% ==================== CHAPTER 5 ====================
\chapter{3DGS 인체 재구성}

\section{정적 배경 3DGS 학습}

동적 인체를 재구성하기 전에, 먼저 빈 장면(인체 없는 상태)의 정적 3DGS를 학습한다.
이 배경 모델은 이후 동적 인체 분리(foreground/background decomposition)에 활용된다.

\begin{pipelinebox}[title={배경 3DGS 학습 단계}]
\begin{enumerate}[nosep]
  \item 인체가 없는 상태에서 다시점 촬영 (또는 인체 마스크로 배경만 추출)
  \item COLMAP으로 초기 점군 생성
  \item 3DGS 학습: 30,000 iterations 기준 약 20--30분 (RTX 4090)
  \item 렌더링하여 배경 재구성 품질 확인
\end{enumerate}
\end{pipelinebox}

3DGS의 각 Gaussian $g_k$는 다음 파라미터로 정의된다:
\begin{equation}
  g_k = \left\{ \vmu_k,\; \mSigma_k,\; \alpha_k,\; \mathbf{c}_k \right\}
\end{equation}
여기서 $\vmu_k \in \RR^3$은 중심, $\mSigma_k \in \RR^{3\times3}$은 공분산(스케일+회전으로 분해),
$\alpha_k \in [0,1]$은 불투명도, $\mathbf{c}_k$는 SH 계수(색상)이다.

\section{동적 인체: GART 학습 파이프라인}

GART(Gaussian Articulated Representation of Transformable humans)는
SMPL-X의 스키닝 가중치를 활용하여 정규 공간(canonical space)의 Gaussian을
각 프레임의 자세로 변환하는 구조이다.

\begin{figure}[H]
\centering
\begin{tikzpicture}[
  node distance=1.0cm and 2.0cm,
  block/.style={rectangle, rounded corners, draw=defblue, fill=lightblue,
    minimum width=2.5cm, minimum height=0.7cm, font=\small, align=center},
  arrow/.style={-Stealth, thick, defblue!60},
  label/.style={font=\tiny\color{gray}, align=center}
]
% Canonical space
\node[block] (canon) {Canonical\\Gaussians};
\node[block, right=2.5cm of canon] (lbs) {LBS\\Skinning};
\node[block, right=2.5cm of lbs] (posed) {Posed\\Gaussians};
\node[block, below=1.2cm of posed] (render) {Differentiable\\Splatting};
\node[block, below=1.2cm of render] (loss) {Photometric\\Loss};

% SMPL-X input
\node[block, fill=lightgreen, draw=biogreen, above=0.8cm of lbs] (smplx) {SMPL-X\\$(\vbeta, \vtheta_t)$};

\draw[arrow] (canon) -- (lbs);
\draw[arrow] (lbs) -- (posed);
\draw[arrow] (posed) -- (render);
\draw[arrow] (render) -- (loss);
\draw[arrow, biogreen!60] (smplx) -- (lbs);

% Feedback
\draw[arrow, darkred, dashed] (loss.west) -- ++(-2,0) |- node[label, left, pos=0.25]{gradient} (canon.south);
\end{tikzpicture}
\caption{GART 학습 루프: 정규 공간 Gaussian을 SMPL-X 자세로 변환 후 렌더링하여 학습}
\end{figure}

\subsection{LBS(Linear Blend Skinning)를 통한 변환}

정규 공간의 Gaussian 중심 $\vmu_k^{\text{can}}$을 자세 $\vtheta_t$로 변환한다:

\begin{equation}
  \vmu_k^{\text{posed}} = \left(\sum_{b=1}^{B} w_{kb}\, \mT_b(\vtheta_t)\right) \vmu_k^{\text{can}}
\end{equation}
여기서 $w_{kb}$는 Gaussian $k$의 본(bone) $b$에 대한 스키닝 가중치,
$\mT_b(\vtheta_t) \in SE(3)$는 본 $b$의 변환 행렬이다.

공분산도 함께 변환된다:
\begin{equation}
  \mSigma_k^{\text{posed}} = \mR_{\text{lbs}} \, \mSigma_k^{\text{can}} \, \mR_{\text{lbs}}^\top
\end{equation}
여기서 $\mR_{\text{lbs}}$는 LBS 변환의 회전 성분이다.

\subsection{학습 가능 파라미터}

GART에서 최적화하는 파라미터는:
\begin{itemize}[nosep]
  \item 정규 공간 Gaussian 속성: $\{\vmu_k^{\text{can}}, \mSigma_k^{\text{can}}, \alpha_k, \mathbf{c}_k\}$
  \item 스키닝 가중치 오프셋: $\Delta w_{kb}$ (SMPL-X 기본 가중치에서의 미세 조정)
  \item (선택) 비강체 변형 네트워크: $\Delta\vmu_k = f_\phi(\vmu_k^{\text{can}}, \vtheta_t)$
\end{itemize}

\section{렌더링 기반 자세 정제}

SMPL-X 피팅 결과에 남아 있는 자세 오차를 렌더링 손실로 정제할 수 있다.
이는 분석-합성(analysis-by-synthesis) 접근법이다:

\begin{equation}
  \vtheta_t^* = \arg\min_{\vtheta_t} \left\| I_{\text{rendered}}(\vtheta_t) - I_{\text{gt}} \right\|_1 + \lambda_{\text{SSIM}} \big(1 - \text{SSIM}(I_{\text{rendered}}, I_{\text{gt}})\big)
\end{equation}

\begin{keyidea}
렌더링 기반 자세 정제는 2D 관절 검출이 실패하는 관절(예: 오클루전된 팔꿈치)에서도
포토메트릭 정보를 활용하여 자세를 교정할 수 있다는 점에서 강력하다.
단, 학습 초반에 Gaussian이 충분히 수렴하지 않은 상태에서 자세를 함께 최적화하면
국소해에 빠질 수 있으므로, 자세 정제는 사전 학습된 Gaussian 위에서 수행한다.
\end{keyidea}

\section{학습 모니터링}

\subsection{정량적 지표}

학습 과정에서 모니터링해야 할 핵심 지표:

\begin{table}[H]
\centering
\caption{3DGS 학습 모니터링 지표}
\begin{tabular}{lccc}
\toprule
\textbf{지표} & \textbf{수식} & \textbf{목표값} & \textbf{방향} \\
\midrule
PSNR & $10\log_{10}(1/\text{MSE})$ & $> 30$\,dB & $\uparrow$ \\
SSIM & 구조적 유사도 & $> 0.95$ & $\uparrow$ \\
LPIPS & 지각적 거리 & $< 0.05$ & $\downarrow$ \\
\# Gaussians & --- & $100\text{K}\sim500\text{K}$ & 안정 \\
\bottomrule
\end{tabular}
\end{table}

\subsection{학습 곡선 해석}

\begin{itemize}[nosep]
  \item \textbf{PSNR이 정체}: 학습률 감소 스케줄 확인, Gaussian densification 주기 조정
  \item \textbf{PSNR 급락}: 자세 파라미터 동시 최적화 시 발생 가능 --- 자세 학습률 별도 관리
  \item \textbf{Gaussian 수 폭증}: pruning 임계값 조정, opacity reset 주기 확인
\end{itemize}

\section{디버깅 가이드}

\begin{table}[H]
\centering
\caption{3DGS 인체 재구성 실패 모드와 해결법}
\begin{tabularx}{\textwidth}{lXX}
\toprule
\textbf{실패 모드} & \textbf{원인} & \textbf{해결} \\
\midrule
유령 사지(ghost limbs) & 자세 오차로 정규 공간에서 잘못된 위치에 Gaussian 생성 & 자세 정제 후 재학습, 시간 정규화 강화 \\
얼룩(floater) & 배경의 Gaussian이 인체에 붙음 & 마스크 기반 배경 분리 강화 \\
흐릿한 텍스처 & Gaussian이 너무 크거나 수가 부족 & densification 임계값 낮춤 \\
색 번짐(color bleeding) & SH 차수가 높아 과적합 & SH 차수를 0--1로 제한하고 점진적 증가 \\
\bottomrule
\end{tabularx}
\end{table}


% ==================== CHAPTER 6 ====================
\chapter{생체역학 지표 산출}

\section{SMPL에서 해부학적 관절 각도로}

SMPL-X의 관절 회전 $\vtheta \in \RR^{3}$ (축-각도 표현)은 직접적으로 임상적 의미를 갖지 않는다.
ISB(International Society of Biomechanics) 권장 분해를 통해
해부학적으로 의미 있는 관절 각도로 변환해야 한다.

\subsection{ISB 분해}

어깨 관절을 예로 들면, 관절 회전 행렬 $\mR_{\text{shoulder}}$를
$Y$-$X'$-$Y''$ 오일러 각으로 분해한다:

\begin{equation}
  \mR_{\text{shoulder}} = \mR_y(\alpha) \, \mR_{x'}(\beta) \, \mR_{y''}(\gamma)
\end{equation}
여기서:
\begin{itemize}[nosep]
  \item $\alpha$: 수평면 거상각(plane of elevation)
  \item $\beta$: 거상각(elevation angle)
  \item $\gamma$: 내/외회전(axial rotation)
\end{itemize}

\begin{biomechbox}
투구 동작에서의 핵심 관절 각도:
\begin{itemize}[nosep]
  \item \textbf{어깨 외회전(MER)}: 투구 중 최대 외회전 각도, 일반적으로 160--180$^\circ$
  \item \textbf{팔꿈치 굴곡}: 가속기 시작 시 약 90--100$^\circ$
  \item \textbf{몸통 회전}: 골반 대비 상체 회전, 최대 40--60$^\circ$ 분리
\end{itemize}
\end{biomechbox}

\begin{codebox}[title={SMPL-X 관절 회전에서 ISB 각도 추출}]
\begin{lstlisting}[language=Python]
from scipy.spatial.transform import Rotation

def smplx_to_isb_shoulder(pose_params, parent_orient):
    """Extract ISB shoulder angles from SMPL-X."""
    # SMPL-X uses axis-angle representation
    R_local = Rotation.from_rotvec(pose_params).as_matrix()
    R_parent = Rotation.from_rotvec(parent_orient).as_matrix()

    # Joint rotation in parent frame
    R_joint = R_parent.T @ R_local

    # YXY Euler decomposition (ISB shoulder)
    r = Rotation.from_matrix(R_joint)
    angles = r.as_euler('YXY', degrees=True)

    return {
        'plane_of_elevation': angles[0],
        'elevation': angles[1],
        'axial_rotation': angles[2]
    }
\end{lstlisting}
\end{codebox}

\section{가상 마커 추출과 OpenSim IK}

SMPL-X 메시 표면에 가상 마커를 정의하여 OpenSim의 역운동학(IK)에 입력할 수 있다.
이는 마커리스 모션캡처와 전통적 생체역학 소프트웨어를 연결하는 핵심 단계이다.

\begin{pipelinebox}[title={가상 마커에서 OpenSim IK까지}]
\begin{enumerate}[nosep]
  \item SMPL-X 메시에서 해부학적 랜드마크에 해당하는 정점(vertex) 인덱스 정의
  \item 프레임별 SMPL-X 출력에서 해당 정점의 3D 좌표 추출
  \item TRC (Track Row Column) 형식으로 변환
  \item OpenSim IK 도구에 입력하여 근골격 모델의 관절 각도 산출
\end{enumerate}
\end{pipelinebox}

\begin{equation}
  \vp_{\text{marker}}^{(m)}(t) = \sum_{i \in \mathcal{N}(m)} \alpha_i \, \vv_i(t)
\end{equation}
여기서 $\mathcal{N}(m)$은 마커 $m$ 주변의 정점 집합, $\alpha_i$는 바리센트릭 가중치이다.

\section{역동역학(Inverse Dynamics)}

관절 각도($\vq$), 각속도($\dot{\vq}$), 각가속도($\ddot{\vq}$)와 외력 데이터로부터
관절 모멘트와 힘을 계산한다.

\begin{equation}
  \bm{\tau} = \mathbf{M}(\vq)\ddot{\vq} + \mathbf{C}(\vq, \dot{\vq})\dot{\vq} + \mathbf{G}(\vq) - \mJ^\top \vF_{\text{ext}}
\end{equation}
여기서 $\mathbf{M}$은 질량 행렬, $\mathbf{C}$는 코리올리/원심력 항,
$\mathbf{G}$는 중력 항, $\vF_{\text{ext}}$는 외력(지면반력 등)이다.

\begin{importantnote}
마커리스 시스템에서 역동역학을 수행할 때 가장 큰 도전은 정확한 외력 측정이다.
Vicon + 지면반력 측정기(force plate)가 있다면 직접 사용하고,
그렇지 않으면 접촉 모델(contact model)이나 잔차 축소 알고리즘(residual reduction)을
적용해야 한다.
\end{importantnote}

\section{투구 이벤트 자동 검출}

투구 동작을 구간별로 분석하기 위해 핵심 이벤트를 자동 검출한다.

\begin{table}[H]
\centering
\caption{투구 이벤트 정의}
\begin{tabularx}{\textwidth}{lXl}
\toprule
\textbf{이벤트} & \textbf{정의} & \textbf{검출 기준} \\
\midrule
FC (Foot Contact) & 앞발이 지면에 착지 & 앞발 수직 속도 $\approx 0$, 높이 최저 \\
MER (Max External Rotation) & 어깨 최대 외회전 & 외회전 각도 극대값 \\
BR (Ball Release) & 공이 손에서 이탈 & 손 가속도 급변, 공 위치 분기 \\
MIR (Max Internal Rotation) & 어깨 최대 내회전 & 감속기 내회전 극대값 \\
\bottomrule
\end{tabularx}
\end{table}

\begin{codebox}[title={이벤트 검출 알고리즘}]
\begin{lstlisting}[language=Python]
from scipy.signal import argrelextrema

def detect_mer(shoulder_ext_rot, fps=240):
    """Detect Maximum External Rotation event."""
    # Find local maxima in external rotation
    maxima = argrelextrema(shoulder_ext_rot, np.greater,
                           order=int(fps*0.05))[0]
    if len(maxima) == 0:
        return None
    # MER is the global maximum during arm cocking
    mer_idx = maxima[np.argmax(shoulder_ext_rot[maxima])]
    return mer_idx

def detect_foot_contact(foot_height, foot_velocity, fps=240):
    """Detect Foot Contact (stride foot)."""
    # Height below threshold AND velocity near zero
    low = foot_height < np.percentile(foot_height, 10)
    slow = np.abs(foot_velocity) < 0.5  # m/s
    candidates = np.where(low & slow)[0]
    return candidates[0] if len(candidates) > 0 else None
\end{lstlisting}
\end{codebox}

\section{핵심 지표 산출}

\subsection{외반 토크 (Valgus Torque)}

팔꿈치 내측측부인대(UCL)에 가해지는 부하의 대리 지표:
\begin{equation}
  \tau_{\text{valgus}} = \tau_{\text{elbow,varus/valgus}} \bigg|_{t = t_{\text{MER}}}
\end{equation}

프로 투수의 경우 일반적으로 60--120\Nm{} 범위이며,
이 값이 높을수록 UCL 손상 위험이 증가한다.

\subsection{최대 외회전 각도}

\begin{equation}
  \text{MER} = \max_{t \in [t_{\text{FC}}, t_{\text{BR}}]} \theta_{\text{ext\_rot}}(t)
\end{equation}

\subsection{팔 속도}

볼 릴리스 시점의 손 선속도:
\begin{equation}
  v_{\text{hand}} = \left\|\frac{d\vp_{\text{hand}}}{dt}\bigg|_{t=t_{\text{BR}}}\right\|
\end{equation}

\begin{injurybox}
\textbf{UCL 부하 기준값}:
\begin{itemize}[nosep]
  \item $\tau_{\text{valgus}} < 60\,\Nm$: 일반적으로 안전
  \item $60 \leq \tau_{\text{valgus}} < 90\,\Nm$: 주의 필요 (반복 부하 누적)
  \item $\tau_{\text{valgus}} \geq 90\,\Nm$: 높은 위험 --- 메카닉 교정 권장
\end{itemize}
이 기준값은 체중, 연령, 투구 유형에 따라 달라지므로 절대적 기준으로 사용해서는 안 된다.
\end{injurybox}


% ==================== CHAPTER 7 ====================
\chapter{Vicon 검증과 오차 분석}

마커리스 파이프라인의 신뢰성을 확립하기 위해 Vicon gold-standard와의 체계적 비교가 필수적이다.

\section{시간 정합 확인}

두 시스템의 데이터를 비교하기 전에 시간축이 정확히 정합되어야 한다.
크로스-코릴레이션(cross-correlation)으로 최적 시간 오프셋을 찾는다:

\begin{equation}
  \Delta t^* = \arg\max_{\Delta t} \sum_{t} s_{\text{vicon}}(t) \cdot s_{\text{markerless}}(t + \Delta t)
\end{equation}

여기서 $s$는 빠르게 변하는 신호(예: 손 가속도의 크기)를 사용한다.

\begin{codebox}[title={크로스-코릴레이션 시간 정합}]
\begin{lstlisting}[language=Python]
from scipy.signal import correlate

def find_time_offset(sig1, sig2, fs):
    """Find time offset between two signals."""
    # Normalize
    s1 = (sig1 - sig1.mean()) / sig1.std()
    s2 = (sig2 - sig2.mean()) / sig2.std()

    corr = correlate(s1, s2, mode='full')
    lag = np.arange(-len(s2)+1, len(s1))
    offset_samples = lag[np.argmax(corr)]
    offset_sec = offset_samples / fs
    return offset_sec
\end{lstlisting}
\end{codebox}

\section{자세 정확도 평가}

\subsection{MPJPE (Mean Per Joint Position Error)}

가장 기본적인 3D 자세 정확도 지표:
\begin{equation}
  \text{MPJPE} = \frac{1}{T \cdot J} \sum_{t=1}^{T}\sum_{j=1}^{J} \left\| \hat{\vp}_j(t) - \vp_j^{\text{gt}}(t) \right\|_2
\end{equation}

Procrustes 정합 후의 PA-MPJPE도 함께 보고한다:
\begin{equation}
  \text{PA-MPJPE} = \frac{1}{TJ} \sum_{t,j} \left\| s^*\mR^*\hat{\vp}_j(t) + \vt^* - \vp_j^{\text{gt}}(t) \right\|_2
\end{equation}

\subsection{관절각 RMSE}

관절별 각도 오차를 평가한다:
\begin{equation}
  \text{RMSE}_{\theta} = \sqrt{\frac{1}{T}\sum_{t=1}^{T} \left(\theta_{\text{markerless}}(t) - \theta_{\text{vicon}}(t)\right)^2}
\end{equation}

\begin{table}[H]
\centering
\caption{기대 정확도 수준 (투구 동작 기준)}
\begin{tabular}{lcc}
\toprule
\textbf{지표} & \textbf{우수} & \textbf{수용 가능} \\
\midrule
MPJPE & $< 20$\,mm & $< 40$\,mm \\
PA-MPJPE & $< 15$\,mm & $< 30$\,mm \\
어깨 각도 RMSE & $< 5^\circ$ & $< 10^\circ$ \\
팔꿈치 각도 RMSE & $< 3^\circ$ & $< 8^\circ$ \\
\bottomrule
\end{tabular}
\end{table}

\section{운동역학 정확도}

토크와 힘의 비교는 역운동학 오차가 미분으로 증폭되므로 더 엄격한 기준이 필요하다.

\begin{equation}
  \text{nRMSE}_\tau = \frac{\text{RMSE}_\tau}{\tau_{\max}^{\text{gt}} - \tau_{\min}^{\text{gt}}} \times 100\,\%
\end{equation}

일반적으로 nRMSE $< 10\%$를 목표로 한다.

\section{Bland--Altman 분석과 ICC}

\subsection{Bland--Altman 분석}

두 측정 방법 간의 일치도를 평가하는 표준 방법이다.

\begin{definition}[Bland--Altman 플롯]
각 측정에 대해:
\begin{itemize}[nosep]
  \item $x$축: 두 방법의 평균 $\bar{x}_i = (x_i^{(1)} + x_i^{(2)}) / 2$
  \item $y$축: 두 방법의 차이 $d_i = x_i^{(1)} - x_i^{(2)}$
\end{itemize}
바이어스(bias)는 $\bar{d} = \frac{1}{n}\sum d_i$이고,
95\% 일치한계(limits of agreement, LoA)는 $\bar{d} \pm 1.96 \cdot s_d$이다.
\end{definition}

\begin{figure}[H]
\centering
\begin{tikzpicture}
\begin{axis}[
  width=10cm, height=6cm,
  xlabel={평균 (Vicon + Markerless) / 2},
  ylabel={차이 (Markerless $-$ Vicon)},
  xmin=100, xmax=200,
  ymin=-15, ymax=15,
  grid=both,
  grid style={gray!20},
  title={Bland--Altman: 어깨 외회전 ($^\circ$)},
  title style={font=\small\bfseries}
]
% Simulated data points
\addplot[only marks, mark=*, mark size=1.5pt, blue!60] coordinates {
  (120,2.1)(125,-1.3)(130,3.5)(135,-0.8)(140,1.2)(145,-2.1)(150,0.5)
  (155,3.2)(160,-1.5)(165,2.8)(170,-0.3)(175,1.9)(180,-2.5)(185,0.7)
  (115,1.8)(128,-0.5)(133,2.3)(138,-1.7)(143,0.9)(148,-2.8)(153,1.1)
  (158,3.0)(163,-0.9)(168,2.1)(173,-1.2)(178,0.4)(183,-2.0)(188,1.5)
  (122,2.5)(127,-1.8)(132,0.3)(137,3.1)(142,-0.6)(147,1.7)(152,-2.3)
  (157,0.8)(162,2.7)(167,-1.4)(172,0.1)(177,3.3)(182,-0.7)(190,1.0)
};
% Bias line
\addplot[red, thick, domain=100:200] {0.6};
\node[red, font=\small] at (axis cs:195,1.5) {bias};
% LoA
\addplot[red, dashed, domain=100:200] {0.6+1.96*1.8};
\addplot[red, dashed, domain=100:200] {0.6-1.96*1.8};
\node[red, font=\tiny] at (axis cs:195,4.8) {+1.96 SD};
\node[red, font=\tiny] at (axis cs:195,-3.6) {$-$1.96 SD};
\end{axis}
\end{tikzpicture}
\caption{Bland--Altman 플롯 예시: 어깨 외회전 각도}
\end{figure}

\subsection{ICC (Intraclass Correlation Coefficient)}

\begin{equation}
  \text{ICC}(3,1) = \frac{\text{MS}_R - \text{MS}_E}{\text{MS}_R + (k-1)\text{MS}_E}
\end{equation}
여기서 $\text{MS}_R$은 피험자 간 평균제곱, $\text{MS}_E$는 잔차 평균제곱, $k$는 측정 방법 수(=2)이다.

\begin{table}[H]
\centering
\caption{ICC 해석 기준 \cite{koo2016}}
\begin{tabular}{ll}
\toprule
\textbf{ICC 값} & \textbf{해석} \\
\midrule
$< 0.50$ & Poor \\
$0.50 \sim 0.75$ & Moderate \\
$0.75 \sim 0.90$ & Good \\
$> 0.90$ & Excellent \\
\bottomrule
\end{tabular}
\end{table}

\section{오차 원인 분석}

파이프라인의 각 단계에서 발생하는 오차를 분리하여 분석한다.

\begin{figure}[H]
\centering
\begin{tikzpicture}[
  node distance=0.8cm,
  block/.style={rectangle, rounded corners, draw=darkred, fill=lightred,
    minimum width=4cm, minimum height=0.6cm, font=\small, align=center},
  arrow/.style={-Stealth, thick, darkred!50}
]
\node[block] (e1) {2D 검출 오차\\(RMSE $\sim$3--5\,px)};
\node[block, below=of e1] (e2) {삼각측량 오차\\(보정 정밀도 의존)};
\node[block, below=of e2] (e3) {SMPL-X 피팅 오차\\(모델 표현력 한계)};
\node[block, below=of e3] (e4) {필터링 왜곡\\(위상 지연, 피크 감쇠)};
\node[block, below=of e4] (e5) {관절각 분해 오차\\(짐벌 락, 좌표계 정의)};
\node[block, below=of e5] (e6) {역동역학 증폭\\(미분으로 인한 노이즈 증폭)};

\draw[arrow] (e1) -- (e2);
\draw[arrow] (e2) -- (e3);
\draw[arrow] (e3) -- (e4);
\draw[arrow] (e4) -- (e5);
\draw[arrow] (e5) -- (e6);

\node[right=1.5cm of e1, font=\small\color{gray}] {$\sim$5\,mm};
\node[right=1.5cm of e2, font=\small\color{gray}] {$\sim$10\,mm};
\node[right=1.5cm of e3, font=\small\color{gray}] {$\sim$15\,mm};
\node[right=1.5cm of e4, font=\small\color{gray}] {$\sim$2$^\circ$};
\node[right=1.5cm of e5, font=\small\color{gray}] {$\sim$5$^\circ$};
\node[right=1.5cm of e6, font=\small\color{gray}] {$\sim$10\,\Nm};
\end{tikzpicture}
\caption{오차 전파 경로와 각 단계의 예상 오차 크기}
\end{figure}

\begin{keyidea}
오차 분석의 핵심 교훈: 역동역학 지표의 정확도는 파이프라인 \textbf{최초 단계}(2D 검출, 카메라 보정)의
정밀도에 의해 근본적으로 제한된다. 마지막 단계(역동역학)만 개선해서는 전체 정확도를 높일 수 없다.
``쓰레기가 들어가면 쓰레기가 나온다(GIGO)'' 원칙이 생체역학 파이프라인에서도 그대로 적용된다.
\end{keyidea}


% ==================== CHAPTER 8 ====================
\chapter{시각화와 보고서 생성}

\section{3DGS 렌더링 + 생체역학 오버레이}

GART로 재구성한 포토리얼리스틱 인체 위에 생체역학 정보를 시각적으로 오버레이하면,
움직임과 부하를 직관적으로 이해할 수 있다.

\subsection{관절 토크 컬러맵}

각 관절의 토크 크기를 컬러맵으로 매핑하여 메시 위에 표시한다:

\begin{equation}
  \text{color}(j, t) = \text{colormap}\!\left(\frac{|\tau_j(t)| - \tau_{\min}}{\tau_{\max} - \tau_{\min}}\right)
\end{equation}

\begin{codebox}[title={토크 오버레이 렌더링}]
\begin{lstlisting}[language=Python]
import matplotlib.cm as cm

def overlay_torque_on_mesh(vertices, faces, joint_torques,
                           joint_indices, skinning_weights):
    """Map joint torques to vertex colors."""
    # Normalize torques to [0, 1]
    t_norm = (np.abs(joint_torques) - t_min) / (t_max - t_min)
    t_norm = np.clip(t_norm, 0, 1)

    # Map to vertices via skinning weights
    vertex_values = skinning_weights @ t_norm

    # Apply colormap (coolwarm: blue=low, red=high)
    colors = cm.coolwarm(vertex_values)[:, :3]
    return colors
\end{lstlisting}
\end{codebox}

\subsection{벡터 필드 시각화}

관절 모멘트의 방향과 크기를 3D 화살표로 표시한다.
특히 팔꿈치 외반 모멘트의 방향은 UCL 부하를 직관적으로 보여준다.

\section{시계열 관절 각도/토크 플롯}

투구 사이클 전체에 걸친 관절 각도와 토크의 시간 변화를 그래프로 나타낸다.

\begin{codebox}[title={시계열 플롯 생성}]
\begin{lstlisting}[language=Python]
import matplotlib.pyplot as plt

fig, axes = plt.subplots(3, 1, figsize=(12, 8), sharex=True)

# Shoulder external rotation
axes[0].plot(time, shoulder_ext_rot, 'b-', label='Markerless')
axes[0].plot(time, vicon_ext_rot, 'r--', label='Vicon')
axes[0].axvline(t_mer, color='gray', ls=':', label='MER')
axes[0].set_ylabel('External Rot. (deg)')
axes[0].legend()

# Elbow flexion
axes[1].plot(time, elbow_flex, 'b-')
axes[1].plot(time, vicon_elbow, 'r--')
axes[1].set_ylabel('Elbow Flex. (deg)')

# Valgus torque
axes[2].plot(time, valgus_torque, 'b-')
axes[2].plot(time, vicon_torque, 'r--')
axes[2].set_ylabel('Valgus Torque (Nm)')
axes[2].set_xlabel('Time (s)')

plt.tight_layout()
plt.savefig('pitching_timeseries.pdf', dpi=300)
\end{lstlisting}
\end{codebox}

\section{투구별 비교 대시보드}

여러 투구를 시간 정규화(time-normalized)하여 중첩 비교한다.
투구 사이클은 FC$\rightarrow$MER$\rightarrow$BR$\rightarrow$MIR의 4개 이벤트로 구분하고,
각 구간을 0--100\%로 정규화한다.

\begin{table}[H]
\centering
\caption{투구별 비교 대시보드 항목}
\begin{tabularx}{\textwidth}{lX}
\toprule
\textbf{항목} & \textbf{내용} \\
\midrule
3D 시각화 & 4개 이벤트 시점의 3DGS 렌더링 스냅샷 \\
각도 프로파일 & 어깨/팔꿈치 각도의 시간 정규화 곡선 (평균$\pm$SD 밴드) \\
토크 프로파일 & 팔꿈치 외반 토크의 시간 정규화 곡선 \\
핵심 지표 표 & MER, 외반 토크, 팔 속도, 보폭 등 스칼라 지표 \\
\bottomrule
\end{tabularx}
\end{table}

\section{자동 보고서 생성 파이프라인}

분석 결과를 자동으로 PDF 보고서로 생성하는 파이프라인을 구축한다.

\begin{codebox}[title={Jinja2 기반 LaTeX 보고서 자동 생성}]
\begin{lstlisting}[language=Python]
from jinja2 import Environment, FileSystemLoader

env = Environment(loader=FileSystemLoader('templates/'))
template = env.get_template('pitching_report.tex.j2')

report = template.render(
    subject_id='S01',
    date='2026-03-15',
    pitches=pitch_results,  # list of per-pitch metrics
    summary_stats=summary,  # aggregate statistics
    figure_paths=fig_paths  # generated plot paths
)

with open('reports/S01_report.tex', 'w') as f:
    f.write(report)

# Compile
os.system('xelatex reports/S01_report.tex')
\end{lstlisting}
\end{codebox}


% ==================== CHAPTER 9 ====================
\chapter{수학적 부록}

\section{Procrustes 정합}

\subsection{문제 정의}

$N$개의 대응 점 쌍 $\{(\va_i, \mathbf{b}_i)\}_{i=1}^{N}$이 주어졌을 때,
최적의 강체 변환(회전 $\mR$, 이동 $\vt$)과 스케일 $s$를 구한다:

\begin{equation}
  \min_{\mR \in SO(3),\, \vt,\, s} \sum_{i=1}^{N} \left\| \mathbf{b}_i - (s\mR\va_i + \vt) \right\|^2
\end{equation}

\subsection{풀이}

\textbf{Step 1.} 중심 이동:
\begin{equation}
  \bar{\va} = \frac{1}{N}\sum_i \va_i, \quad \bar{\mathbf{b}} = \frac{1}{N}\sum_i \mathbf{b}_i, \quad
  \tilde{\va}_i = \va_i - \bar{\va}, \quad \tilde{\mathbf{b}}_i = \mathbf{b}_i - \bar{\mathbf{b}}
\end{equation}

\textbf{Step 2.} 교차공분산 행렬과 SVD:
\begin{equation}
  \mH = \sum_{i=1}^{N} \tilde{\va}_i \tilde{\mathbf{b}}_i^\top, \quad
  \mH = \mathbf{U}\bm{\Lambda}\mathbf{V}^\top
\end{equation}

\textbf{Step 3.} 최적 회전:
\begin{equation}
  \mR^* = \mathbf{V} \begin{pmatrix} 1 & & \\ & 1 & \\ & & \det(\mathbf{V}\mathbf{U}^\top) \end{pmatrix} \mathbf{U}^\top
\end{equation}

\textbf{Step 4.} 최적 스케일과 이동:
\begin{equation}
  s^* = \frac{\text{tr}(\bm{\Lambda} \cdot \text{diag}(1,1,\det(\mathbf{V}\mathbf{U}^\top)))}{\sum_i \|\tilde{\va}_i\|^2}, \quad
  \vt^* = \bar{\mathbf{b}} - s^*\mR^*\bar{\va}
\end{equation}

\section{Butterworth 필터 설계}

\subsection{아날로그 프로토타입}

$n$차 Butterworth 저역통과 필터의 전달 함수:
\begin{equation}
  H(s) = \frac{1}{\prod_{k=1}^{n} (s - p_k)}, \quad
  p_k = e^{j\pi(2k+n-1)/(2n)}, \quad k = 1, \ldots, n
\end{equation}

\subsection{디지털 변환}

쌍일차 변환(bilinear transform)으로 아날로그 필터를 디지털로 변환한다:
\begin{equation}
  s = \frac{2}{T_s} \cdot \frac{1 - z^{-1}}{1 + z^{-1}}
\end{equation}

사전 왜곡(pre-warping)을 적용하여 차단 주파수의 정확성을 보장한다:
\begin{equation}
  \Omega_c = \frac{2}{T_s} \tan\!\left(\frac{\omega_c T_s}{2}\right)
\end{equation}

\subsection{양방향 필터링(Zero-Phase)}

\texttt{filtfilt}은 필터를 순방향$\cdot$역방향으로 두 번 적용하여 위상 지연을 제거한다:
\begin{equation}
  |H_{\text{zp}}(f)|^2 = |H(f)|^4 = \frac{1}{\left(1 + (f/f_c)^{2n}\right)^2}
\end{equation}
유효 차수가 $2n$이 되므로, 실제 차단 특성은 단방향 대비 두 배 급격해진다.

\section{Bland--Altman 분석의 통계적 기초}

\subsection{가정}

\begin{itemize}[nosep]
  \item 차이 $d_i$가 정규분포를 따른다: $d_i \sim \mathcal{N}(\mu_d, \sigma_d^2)$
  \item 차이의 크기가 측정값의 크기에 의존하지 않는다 (비례 바이어스 없음)
  \item 관측치가 서로 독립이다
\end{itemize}

\subsection{일치한계(Limits of Agreement)}

\begin{equation}
  \text{LoA} = \bar{d} \pm 1.96 \cdot s_d
\end{equation}
여기서 $\bar{d} = \frac{1}{n}\sum_{i=1}^{n} d_i$, $s_d = \sqrt{\frac{1}{n-1}\sum_{i=1}^{n}(d_i - \bar{d})^2}$이다.

95\% 신뢰구간을 포함한 LoA:
\begin{equation}
  \text{SE}(\text{LoA}) = \sqrt{\frac{1}{n} + \frac{1.96^2}{2(n-1)}} \cdot s_d
\end{equation}

\subsection{비례 바이어스 검정}

차이 $d_i$를 평균 $\bar{x}_i$에 대해 회귀 분석하여 기울기가 0인지 검정한다:
\begin{equation}
  d_i = \beta_0 + \beta_1 \bar{x}_i + \epsilon_i
\end{equation}
$\beta_1$이 유의하면 비례 바이어스가 존재하며, 이 경우 로그 변환 후 재분석하거나
회귀 기반 LoA를 사용한다.

\section{ICC 계산}

\subsection{Two-Way Mixed Model, Absolute Agreement}

ICC(3,1)의 분산분석 모형:
\begin{equation}
  x_{ij} = \mu + r_i + c_j + e_{ij}
\end{equation}
여기서 $r_i$는 피험자 효과(random), $c_j$는 방법 효과(fixed), $e_{ij}$는 잔차이다.

\begin{equation}
  \text{ICC}(3,1) = \frac{\text{MS}_R - \text{MS}_E}{\text{MS}_R + (k-1)\text{MS}_E}
\end{equation}

\subsection{신뢰구간}

ICC의 95\% 신뢰구간은 $F$-분포를 이용하여 계산한다:

\begin{equation}
\begin{aligned}
  F_L &= \frac{F_0}{F_{1-\alpha/2, n-1, (n-1)(k-1)}}, \quad
  F_U = F_0 \cdot F_{1-\alpha/2, (n-1)(k-1), n-1} \\
  \text{ICC}_L &= \frac{F_L - 1}{F_L + k - 1}, \quad
  \text{ICC}_U = \frac{F_U - 1}{F_U + k - 1}
\end{aligned}
\end{equation}
여기서 $F_0 = \text{MS}_R / \text{MS}_E$이다.

\begin{codebox}[title={ICC 계산 구현}]
\begin{lstlisting}[language=Python]
import pingouin as pg
import pandas as pd

def compute_icc(vicon_data, markerless_data, metric_name):
    """Compute ICC(3,1) for a single metric."""
    n = len(vicon_data)
    df = pd.DataFrame({
        'subject': list(range(n)) * 2,
        'method': ['vicon']*n + ['markerless']*n,
        'value': np.concatenate([vicon_data, markerless_data])
    })

    icc = pg.intraclass_corr(
        data=df, targets='subject',
        raters='method', ratings='value'
    )
    # ICC3 = two-way mixed, single measures
    result = icc.set_index('Type').loc['ICC3']
    return result['ICC'], result['CI95%']
\end{lstlisting}
\end{codebox}


% ==================== REFERENCES ====================
\begin{thebibliography}{99}

\bibitem{winter2009}
Winter, D.A. (2009). \textit{Biomechanics and Motor Control of Human Movement}. 4th ed. Wiley.

\bibitem{zhang2000}
Zhang, Z. (2000). A flexible new technique for camera calibration. \textit{IEEE Trans. PAMI}, 22(11), 1330--1334.

\bibitem{xu2022vitpose}
Xu, Y., Zhang, J., Zhang, Q., \& Tao, D. (2022). ViTPose: Simple Vision Transformer Baselines for Human Pose Estimation. \textit{NeurIPS}.

\bibitem{jiang2023rtmpose}
Jiang, T., Lu, P., Zhang, L., et al. (2023). RTMPose: Real-Time Multi-Person Pose Estimation based on MMPose. \textit{arXiv:2303.07399}.

\bibitem{easymocap}
Dong, Z., Song, J., Chen, X., et al. (2021). EasyMocap: Open source toolbox for markerless multi-view motion capture.

\bibitem{pavlakos2019}
Pavlakos, G., Choutas, V., Ghorbani, N., et al. (2019). Expressive body capture: 3D hands, face, and body from a single image. \textit{CVPR}.

\bibitem{kerbl2023}
Kerbl, B., Kopanas, G., Leimk\"uhler, T., \& Drettakis, G. (2023). 3D Gaussian Splatting for Real-Time Radiance Field Rendering. \textit{ACM TOG (SIGGRAPH)}.

\bibitem{lei2024gart}
Lei, J., Wang, Y., Pavlakos, G., \& Daniilidis, K. (2024). GART: Gaussian Articulated Template Models. \textit{CVPR}.

\bibitem{koo2016}
Koo, T.K. \& Li, M.Y. (2016). A guideline of selecting and reporting intraclass correlation coefficients for reliability research. \textit{J Chiropr Med}, 15(2), 155--163.

\bibitem{fleisig2011}
Fleisig, G.S., Barrentine, S.W., Zheng, N., et al. (2011). Kinetic comparison among the fastball, curveball, change-up, and slider in collegiate baseball pitchers. \textit{Am J Sports Med}.

\bibitem{schoenberger2016sfm}
Sch\"onberger, J.L. \& Frahm, J.-M. (2016). Structure-from-Motion Revisited. \textit{CVPR}.

\bibitem{bland1986}
Bland, J.M. \& Altman, D.G. (1986). Statistical methods for assessing agreement between two methods of clinical measurement. \textit{The Lancet}, 327(8476), 307--310.

\bibitem{wu2005isb}
Wu, G., van der Helm, F.C.T., et al. (2005). ISB recommendation on definitions of joint coordinate systems. \textit{J Biomech}, 38(5), 981--992.

\bibitem{loper2015smpl}
Loper, M., Mahmood, N., Romero, J., Pons-Moll, G., \& Black, M.J. (2015). SMPL: A skinned multi-person linear model. \textit{ACM TOG (SIGGRAPH Asia)}.

\bibitem{pavlakos2019smplx}
Pavlakos, G., Choutas, V., Ghorbani, N., et al. (2019). Expressive body capture: 3D hands, face, and body from a single image. \textit{CVPR}.

\end{thebibliography}


% ==================== GLOSSARY ====================
\chapter*{용어집}
\addcontentsline{toc}{chapter}{용어집}

\begin{longtable}{p{4cm}p{10cm}}
\toprule
\textbf{용어} & \textbf{설명} \\
\midrule
\endhead

3DGS & 3D Gaussian Splatting. 3D 공간의 Gaussian 원소를 래스터화하여 장면을 렌더링하는 기법 \\
BA & Bundle Adjustment. 카메라 자세와 3D 점 위치를 동시에 최적화하는 비선형 최소제곱 \\
Bland--Altman & 두 측정 방법 간의 일치도를 평가하는 통계적 방법 \\
COLMAP & Structure from Motion과 Multi-View Stereo를 수행하는 오픈소스 도구 \\
DLT & Direct Linear Transform. 다시점 삼각측량의 선형 해법 \\
FC & Foot Contact. 투구 시 앞발의 지면 접촉 이벤트 \\
GART & Gaussian Articulated Representation of Transformable humans \\
ICC & Intraclass Correlation Coefficient. 측정 신뢰도 지표 \\
ID & Inverse Dynamics. 운동학 데이터로부터 관절 힘/모멘트를 역으로 계산 \\
IK & Inverse Kinematics. 관절 각도를 역으로 계산하는 방법 \\
ISB & International Society of Biomechanics. 관절 좌표계의 표준을 정의 \\
LBS & Linear Blend Skinning. 스켈레톤 기반 메시 변형 기법 \\
LoA & Limits of Agreement. Bland--Altman 분석의 95\% 일치한계 \\
LPIPS & Learned Perceptual Image Patch Similarity. 지각적 이미지 유사도 지표 \\
MER & Maximum External Rotation. 투구 중 어깨의 최대 외회전 \\
MIR & Maximum Internal Rotation. 투구 중 어깨의 최대 내회전 \\
MPJPE & Mean Per Joint Position Error. 3D 자세 추정 오차 지표 \\
MVS & Multi-View Stereo. 다시점 이미지에서 밀집 3D 복원 \\
PA-MPJPE & Procrustes Aligned MPJPE. 최적 강체 정합 후의 관절 위치 오차 \\
PSNR & Peak Signal-to-Noise Ratio. 이미지 품질 지표 \\
RANSAC & Random Sample Consensus. 이상치에 강건한 모델 추정 \\
SfM & Structure from Motion. 다시점 이미지에서 카메라 자세와 희소 3D 복원 \\
SH & Spherical Harmonics. 방향 의존적 색상을 표현하는 기저 함수 \\
SMPL-X & Skinned Multi-Person Linear model eXpressive. 손과 얼굴을 포함하는 파라메트릭 인체 모델 \\
SSIM & Structural Similarity Index. 구조적 이미지 유사도 지표 \\
UCL & Ulnar Collateral Ligament. 팔꿈치 내측측부인대 \\
VPoser & Variational Human Pose Prior. 학습된 자세 사전 확률 모델 \\
\bottomrule
\end{longtable}

\end{document}
